% ============================================================
%  ANSWER KEY — Practices 5, 6, 7, 8
%  (Pure content, no preamble / \begin{document})
% ============================================================

% ============================================================
\section*{Practice 5: Combinatorics --- Basic Rules}
% ============================================================

% ------------------------------------------------------------
%  Section 1: Fundamental Rules (Sum and Product Rules)
% ------------------------------------------------------------
\subsection*{Section 1: Sum and Product Rules}

\subsubsection*{Problem 1}
Count of integers from 0 to 1000 containing at least one digit 6.

Complement: count integers with \emph{no} digit 6.
\begin{itemize}
  \item 1-digit numbers (0--9): 9 have no 6 (all except 6).
  \item 2-digit numbers (10--99): first digit from $\{1,\dots,9\}\setminus\{6\}=8$, second digit from $\{0,\dots,9\}\setminus\{6\}=9$. Total $8\cdot9=72$.
  \item 3-digit numbers (100--999): $8\cdot9\cdot9=648$.
  \item The number 1000 contains no 6: 1 more.
\end{itemize}
No-six total: $9+72+648+1=730$.
All integers from 0 to 1000: $1001$.

$\boxed{1001 - 730 = 271}$

\subsubsection*{Problem 2}
Same problem as Problem 1 (the exercise sheet notes the apparent duplication).

$\boxed{271}$

\subsubsection*{Problem 3}
Choose two books of different subjects from 15 Informatics, 12 Math, 10 Physics.

By the sum/product rules:
$15\cdot12 + 15\cdot10 + 12\cdot10 = 180+150+120 = \boxed{450}$.

\subsubsection*{Problem 4}
Odd numbers between 100 and 1000 (i.e.\ 3-digit odd numbers 101, 103, \ldots, 999).

Last digit: 5 choices ($\{1,3,5,7,9\}$).
First digit: 9 choices ($1$--$9$).
Middle digit: 10 choices ($0$--$9$).

$\boxed{9 \times 10 \times 5 = 450}$

\subsubsection*{Problem 5}
Natural numbers less than 700 ($\{1,\dots,699\}$):

(a) Divisible by 5: $\lfloor 699/5\rfloor = \boxed{139}$.

(b) Divisible by 3: $\lfloor 699/3\rfloor = \boxed{233}$.

(c) Divisible by both 3 and 5 (i.e.\ by 15): $\lfloor 699/15\rfloor = \boxed{46}$.

\subsubsection*{Problem 6}
Yerevan $\to$ Moscow: 15 flights. Moscow $\to$ London: 20 flights.
By the product rule: $\boxed{15 \times 20 = 300}$.

\subsubsection*{Problem 7}
4-digit numbers with at least two identical digits.

Total 4-digit numbers: $9\times 10^3 = 9000$.
4-digit numbers with all digits distinct: first digit 9 choices, second 9, third 8, fourth 7, giving $9\times 9\times 8\times 7 = 4536$.

$\boxed{9000 - 4536 = 4464}$

\subsubsection*{Problem 8}
7-digit numbers where first and last digits are different.

Total 7-digit numbers: $9\times 10^6$.
7-digit numbers where first = last digit: first digit 9 choices, last digit fixed equal to first (1 choice), middle 5 digits $10^5$ each. Total same: $9\times 10^5$.

Answer: $9\times 10^6 - 9\times 10^5 = 9\times 10^5(10-1)=9\times 10^5\times 9$.

$\boxed{9\times 10^6 - 9\times 10^5 = 8{,}100{,}000}$

\subsubsection*{Problem 9}
6-digit numbers where adjacent digits are different (digits 0--9, leading digit $\neq 0$).

Position 1: 9 choices (1--9). Each subsequent position: 9 choices (any digit except the preceding one).

$\boxed{9 \times 9^5 = 9^6 = 531{,}441}$

\subsubsection*{Problem 10}
7-digit numbers: even positions (2,4,6) hold odd digits; odd positions (1,3,5,7) hold even digits.

Positions 2,4,6: 5 odd digits each $\Rightarrow 5^3$.
Position 1 (leading, must be even and $\neq 0$): 4 choices $\{2,4,6,8\}$.
Positions 3,5,7: 5 even digits each $\Rightarrow 5^3$.

$\boxed{4 \times 5^3 \times 5^3 = 4 \times 15{,}625 = 62{,}500}$

% ------------------------------------------------------------
%  Section 2: Permutations and Arrangements
% ------------------------------------------------------------
\subsection*{Section 2: Permutations and Arrangements}

\subsubsection*{Problem 1}
Mappings from a 6-element set to a 3-element set: each of the 6 elements maps to one of 3 images.

$\boxed{3^6 = 729}$

\subsubsection*{Problem 2}
Injective mappings from a 6-element set to a 3-element set.
Since $6 > 3$, by the Pigeonhole Principle no injection exists.

$\boxed{0}$

\subsubsection*{Problem 3}
5 boys and 6 girls in a line; boys occupy the first 5 positions.

Arrange 5 boys in the first 5 positions: $5!$.
Arrange 6 girls in the remaining 6 positions: $6!$.

$\boxed{5! \times 6! = 120 \times 720 = 86{,}400}$

\subsubsection*{Problem 4}
6 boys and 7 girls ($13$ people) in a line; first position is a boy.

Choose the boy for position 1: 6 ways. Arrange remaining 12 people: $12!$.

$\boxed{6 \times 12! = 6 \times 479{,}001{,}600 = 2{,}874{,}009{,}600}$

\subsubsection*{Problem 5}
6 boys and 7 girls in a line; first and last positions are boys.

First position: 6 choices. Last position: 5 remaining boys. Middle 11 positions: $11!$.

$\boxed{6 \times 5 \times 11! = 30 \times 39{,}916{,}800 = 1{,}197{,}504{,}000}$

\subsubsection*{Problem 6}
5 boys and 6 girls in a line; boys occupy 5 \emph{consecutive} positions.

Treat the block of 5 boys as a single unit. Together with 6 girls, we have 7 objects to arrange in a line: $7!$ ways. Within the block the 5 boys can be permuted: $5!$ ways.

$\boxed{7! \times 5! = 5040 \times 120 = 604{,}800}$

\subsubsection*{Problem 7}
Arrange 9 distinct people in a line.

$\boxed{9! = 362{,}880}$

\subsubsection*{Problem 8}
8-digit numbers formed from the digits of 23122313.

Digits: 1 appears 2 times, 2 appears 3 times, 3 appears 2 times (wait, let us recount: 2,3,1,2,2,3,1,3).
Digits: 1 $\times 2$, 2 $\times 3$, 3 $\times 3$. Total 8 digits.

Permutations with repetition: $\dfrac{8!}{2!\,3!\,3!} = \dfrac{40320}{2\cdot6\cdot6} = \dfrac{40320}{72}$.

$\boxed{\frac{8!}{2!\,3!\,3!} = 560}$

\subsubsection*{Problem 9}
Words of length 6 from alphabet $\{a, b, g, d, e, z\}$ (vowels: $a, e$; consonants: $b, g, d, z$) where two vowels are adjacent.

Total words of length 6: $6^6 = 46656$.

We count words where the two vowels $a, e$ appear adjacently. Reinterpretation: we form 6-letter words (with repetition) from the 6-letter alphabet, containing at least one pair of adjacent vowels.

Actually, since the problem says ``where two vowels are adjacent'' (likely meaning at least two adjacent vowels), we use complementary counting.

Strings with NO two adjacent vowels: at each position, choose any of 6 letters, but no two consecutive positions are both vowels (from $\{a,e\}$).

Let $f(n)$ = number of such strings. At each position the letter is a vowel (2 choices) or consonant (4 choices).
Recurrence: $f(n) = 4\cdot f(n-1) + 2\cdot 4\cdot f(n-2)$ ...

More carefully: Let $V_n$ = strings of length $n$ ending in a vowel with no two adjacent vowels, $C_n$ = ending in a consonant. Then $V_n = 2\cdot C_{n-1}$, $C_n = 4(V_{n-1}+C_{n-1})$. With $f(n)=V_n+C_n$.

$V_1=2, C_1=4, f(1)=6$.
$V_2=2\cdot4=8, C_2=4(2+4)=24, f(2)=32$.
$V_3=2\cdot24=48, C_3=4(8+24)=128, f(3)=176$.
$V_4=2\cdot128=256, C_4=4(48+128)=704, f(4)=960$.
$V_5=2\cdot704=1408, C_5=4(256+704)=3840, f(5)=5248$.
$V_6=2\cdot3840=7680, C_6=4(1408+3840)=20992, f(6)=28672$.

Answer: $6^6 - 28672 = 46656 - 28672$.

$\boxed{46656 - 28672 = 17984}$

\subsubsection*{Problem 10}
\textbf{Prove:} In any group of 6 people, there are either 3 mutual friends or 3 mutual strangers ($R(3,3)=6$).

\textit{Proof sketch.}
Model the 6 people as vertices of $K_6$ with edges coloured red (friend) or blue (stranger). Pick any vertex $A$; it has 5 edges, so by the Pigeonhole Principle at least 3 share one colour, say red, to vertices $B,C,D$. If any edge among $B,C,D$ is red, it forms a red triangle with $A$; otherwise $B,C,D$ form a blue triangle. \qed

\subsubsection*{Problem 11}
How many numbers must be selected from $\{1,2,\dots,210\}$ to guarantee at least one even number?

There are 105 odd numbers in $\{1,\dots,210\}$. In the worst case we pick all 105 odd numbers first. One more pick guarantees an even number.

$\boxed{106}$

\subsubsection*{Problem 12}
\textbf{Prove:} Among the first 2014 terms of $7, 77, 777, \ldots$, one is divisible by 2013.

\textit{Proof sketch.}
Let $a_k = \underbrace{77\ldots7}_{k}$. Consider the 2014 remainders $a_k \bmod 2013$. If some remainder is 0, done. Otherwise 2014 remainders fall into $\{1,\dots,2012\}$ (2012 classes), so by Pigeonhole two terms $a_i, a_j$ ($i<j$) share a remainder. Then $2013 \mid a_j - a_i = a_{j-i}\cdot 10^i$. Since $\gcd(10,2013)=1$, we get $2013\mid a_{j-i}$, and $1\le j-i\le 2013$. \qed

% ============================================================
\section*{Practice 6: Combinations and Binomial Coefficients}
% ============================================================

\subsubsection*{Problem 1}
Bit strings of length 9 with exactly 5 zeros (and 4 ones).

Choose positions for the 5 zeros: $\binom{9}{5}$.

$\boxed{\binom{9}{5} = 126}$

\subsubsection*{Problem 2}
Committee of 6 teachers and 8 students from 12 teachers and 20 students.

$\boxed{\binom{12}{6}\cdot\binom{20}{8} = 924 \times 125970 = 116{,}396{,}280}$

\subsubsection*{Problem 3}
Choose 13 cards from 52 such that they include 5 cards of the same suit.

This asks for the number of 13-card hands containing at least one suit with 5 or more cards. By inclusion--exclusion on the 4 suits, or noting it is a standard bridge-hand constraint: this is a complex computation. A cleaner reading: choose 13 cards that include \emph{exactly} one group of 5 of the same suit, i.e., at least 5 from some suit.

Using complement (no suit has $\ge 5$ cards, meaning each suit contributes $\le 4$ cards):
\[
N_{\le 4} = \sum_{\substack{a+b+c+d=13\\0\le a,b,c,d\le 4}} \binom{13}{a}\binom{13}{b}\binom{13}{c}\binom{13}{d}
\]
where each suit has 13 cards. The possible distributions $(a,b,c,d)$ with each $\le 4$ summing to 13 are: $(4,4,4,1)$ in $\binom{4}{1}\cdot 4 = 4$ orderings, $(4,4,3,2)$ in $4!/2!=12$, $(4,3,3,3)$ in $4$ orderings.

\begin{align*}
(4,4,4,1):&\quad 4\cdot\binom{13}{4}^3\binom{13}{1}= 4\cdot 715^3\cdot 13 = 4\cdot 365{,}525{,}875\cdot 13 = 19{,}007{,}345{,}500\\
(4,4,3,2):&\quad 12\cdot\binom{13}{4}^2\binom{13}{3}\binom{13}{2}=12\cdot 715^2\cdot 286\cdot 78 = 12\cdot 511225\cdot 22308 = 136{,}852{,}887{,}600\\[2pt]
(4,3,3,3):&\quad 4\cdot\binom{13}{4}\binom{13}{3}^3 = 4\cdot 715\cdot 286^3 = 4\cdot 715\cdot 23{,}393{,}656 = 66{,}905{,}856{,}160
\end{align*}

$N_{\le 4} = 19{,}007{,}345{,}500 + 136{,}852{,}887{,}600 + 66{,}905{,}856{,}160 = 222{,}766{,}089{,}260$.

Total 13-card hands: $\binom{52}{13} = 635{,}013{,}559{,}600$.

$\boxed{635{,}013{,}559{,}600 - 222{,}766{,}089{,}260 = 412{,}247{,}470{,}340}$

\subsubsection*{Problem 4}
Divide 10 people into two teams of 5.

Since the two teams are unordered (indistinguishable labels):
$\dfrac{1}{2}\binom{10}{5} = \dfrac{252}{2}$.

$\boxed{126}$

\subsubsection*{Problem 5}
$A=\{a,b,c,d,e,f,g,h\}$. Number of 5-element subsets containing element $a$.

Fix $a$ in the subset; choose 4 more from the remaining 7 elements.

$\boxed{\binom{7}{4} = 35}$

\subsubsection*{Problem 6}
Coefficient of $x^6 y^4$ in $(2x+3y)^{10}$.

By the binomial theorem: $\binom{10}{6}(2x)^6(3y)^4 = \binom{10}{6}\cdot 2^6\cdot 3^4\cdot x^6 y^4$.

$\binom{10}{6}=210$, $2^6=64$, $3^4=81$.

$\boxed{210\cdot 64\cdot 81 = 1{,}088{,}640}$

\subsubsection*{Problem 7}
Rook at a1 reaching e5, moving only right or up.

From a1 to e5: 4 steps right ($a\to b\to c\to d\to e$) and 4 steps up ($1\to2\to3\to4\to5$). Total 8 steps, choose 4 for right.

$\boxed{\binom{8}{4} = 70}$

\subsubsection*{Problem 8}
Lattice paths from $O(0,0)$ to $A(7,8)$ through $(3,4)$.

$O\to P(3,4)$: $\binom{7}{3}=35$ paths. $P\to A$: $\binom{8}{4}=70$ paths.

$\boxed{35\times 70 = 2450}$

\subsubsection*{Problem 9}
Number of ways to partition 6 into three summands.

\textbf{Ordered} (compositions, parts $\ge 1$): $\binom{5}{2}=\boxed{10}$.

\textbf{Unordered} (integer partitions into 3 parts): $4+1+1$, $3+2+1$, $2+2+2$ $= \boxed{3}$.

\subsubsection*{Problem 10}
Buy 7 cakes from 3 types (combinations with repetition).

$\boxed{\binom{7+3-1}{3-1} = \binom{9}{2} = 36}$

\subsubsection*{Problem 11}
Rearrangements of ``matematika'' (10 letters: a$\times 3$, m$\times 2$, t$\times 2$, e$\times 1$, i$\times 1$, k$\times 1$).

$\boxed{\frac{10!}{3!\cdot 2!\cdot 2!} = \frac{3{,}628{,}800}{24} = 151{,}200}$

\subsubsection*{Problem 12}
Arrange 15 math, 16 informatics, 12 physics books on a shelf (books of the same subject are indistinguishable). Total $15+16+12=43$ books.

$\boxed{\frac{43!}{15!\cdot 16!\cdot 12!} = \binom{43}{15,\,16,\,12}}$

\subsubsection*{Problem 13}
Coefficient of $x_1^4 x_2^3 x_3^2$ in $(x_1+x_2+x_3)^9$.

By the multinomial theorem: $\dfrac{9!}{4!\,3!\,2!}$.

$\boxed{\frac{9!}{4!\,3!\,2!} = \frac{362880}{24\cdot 6\cdot 2} = 1260}$

% ============================================================
\section*{Practice 7: Partitions \& Inclusion-Exclusion}
% ============================================================

\subsubsection*{Problem 1}
Partition a 5-element set into 3 non-empty classes: Stirling number $S(5,3)$.

$S(5,3) = S(4,2)+3\cdot S(4,3) = 7 + 3\cdot 6 = \boxed{25}$.

\subsubsection*{Problem 2}
Distribute $m$ distinct balls into $n$ distinct boxes, no box empty (surjections $m\to n$).

By inclusion--exclusion:
\[
  \boxed{n!\,S(m,n) = \sum_{k=0}^{n}(-1)^k\binom{n}{k}(n-k)^m}
\]

\subsubsection*{Problem 3}
Numbers in $\{1,\dots,1000\}$ not divisible by 2, 3, or 5.

Let $A_d = $ multiples of $d$. By inclusion--exclusion:
\begin{align*}
|A_2|&=500,\; |A_3|=333,\; |A_5|=200,\\
|A_6|&=166,\; |A_{10}|=100,\; |A_{15}|=66,\; |A_{30}|=33.
\end{align*}
$|A_2\cup A_3\cup A_5|=500+333+200-166-100-66+33=734$.

$\boxed{1000-734=266}$

\subsubsection*{Problem 4}
4-digit numbers containing at least one digit 7.

Total 4-digit numbers: 9000. Those with no 7: first digit $8$ choices, each other digit $9$ choices $= 8\times 9^3 = 5832$.

$\boxed{9000-5832=3168}$

\subsubsection*{Problem 5}
Integers in $\{1,\dots,16500\}$ not divisible by 5.

Divisible by 5: $\lfloor 16500/5\rfloor = 3300$.

$\boxed{16500-3300=13200}$

\subsubsection*{Problem 6}
Integers in $\{1,\dots,16500\}$ divisible by neither 5 nor 3.

$|A_3|=5500$, $|A_5|=3300$, $|A_{15}|=1100$.
$|A_3\cup A_5|=5500+3300-1100=7700$.

$\boxed{16500-7700=8800}$

\subsubsection*{Problem 7}
Integers in $\{1,\dots,16500\}$: divisible by 11, not by 5, not by 3.

Universe $S$: multiples of 11, $|S|=1500$. Let $A=$ also div.\ by 3 (div.\ by 33), $B=$ also div.\ by 5 (div.\ by 55).
$|A|=500$, $|B|=300$, $|A\cap B|=\lfloor 16500/165\rfloor=100$.
$|A\cup B|=700$.

$\boxed{1500-700=800}$

\subsubsection*{Problem 8}
Integers in $\{1,\dots,16500\}$ not divisible by 5, 3, or 11.

$|A_3|=5500,\;|A_5|=3300,\;|A_{11}|=1500$.
$|A_{15}|=1100,\;|A_{33}|=500,\;|A_{55}|=300$.
$|A_{165}|=100$.

$|A_3\cup A_5\cup A_{11}|=5500+3300+1500-1100-500-300+100=8500$.

$\boxed{16500-8500=8000}$

\subsubsection*{Problem 9}
30 students; 20 know English, 12 French, 6 both.

$|E\cup F|=20+12-6=26$. Know neither:

$\boxed{30-26=4}$

\subsubsection*{Problem 10}
Triangles in a triangular grid with $n=8$.

Upward-pointing (side $s$, $1\le s\le 8$): $\sum_{s=1}^{8}\binom{10-s}{2} = 36+28+21+15+10+6+3+1=120$.

Downward-pointing (side $s$, $1\le s\le 4$): $\sum_{s=1}^{4}\binom{10-2s}{2}=28+15+6+1=50$.

$\boxed{120+50=170}$

\subsubsection*{Problem 11}
Three carpets of area 3~m$^2$ each in a room of area 6~m$^2$. Prove two overlap by $\ge \frac{1}{2}$~m$^2$.

\textit{Proof sketch.}
Total carpet area $= 9$~m$^2$, room area $= 6$~m$^2$, so the total overlap (counted with multiplicity) is $\ge 9-6=3$~m$^2$. There are $\binom{3}{2}=3$ pairs of carpets. By the Pigeonhole Principle, at least one pair overlaps by $\ge 3/3 = 1$~m$^2 \ge \frac{1}{2}$~m$^2$. \qed

\subsubsection*{Problem 12}
Integer solutions of $x_1+x_2+x_3=40$ with $4\le x_1\le 15$, $9\le x_2\le 18$, $5\le x_3\le 16$.

Substitute $y_i = x_i - \text{lower bound}$: $y_1+y_2+y_3=22$ with $0\le y_1\le 11$, $0\le y_2\le 9$, $0\le y_3\le 11$.

Unrestricted: $\binom{24}{2}=276$.
Violations by inclusion--exclusion: $|A_1|+|A_2|+|A_3|-(|A_1\cap A_2|+|A_1\cap A_3|+|A_2\cap A_3|)+|A_1\cap A_2\cap A_3| = (66+91+66)-(1+0+1)+0 = 221$.

$\boxed{276-221=55}$

% ============================================================
\section*{Practice 8: Recurrence Relations}
% ============================================================

\subsubsection*{Problem 1}
$a_1=3$, $a_2=9$, $a_{n+2}=a_{n+1}-a_n$.

The sequence is periodic with period 6:
$3,\;9,\;6,\;-3,\;-9,\;-6,\;3,\;9,\;\ldots$

Since $200 = 33\times 6 + 2$, we have $a_{200}=a_2$.

$\boxed{a_{200}=9}$

\subsubsection*{Problem 2}
$a_n=5a_{n-1}-6a_{n-2}$, $a_1=2$, $a_2=10$.

Characteristic equation: $r^2-5r+6=0 \Rightarrow r=2,3$.
General solution: $a_n=\alpha\cdot 2^n+\beta\cdot 3^n$.
From initial conditions: $\alpha=-2$, $\beta=2$.

$\boxed{a_n = 2(3^n - 2^n)}$

\subsubsection*{Problem 3}
$a_n=4a_{n-1}-4a_{n-2}$, $a_0=2$, $a_1=6$.

Characteristic equation: $r^2-4r+4=0\Rightarrow r=2$ (double root).
General solution: $a_n=(\alpha+\beta n)\cdot 2^n$.
$a_0=\alpha=2$, $a_1=(\alpha+\beta)\cdot 2 = 2(2+\beta)=6 \Rightarrow \beta=1$.

$\boxed{a_n = (2+n)\cdot 2^n}$

\subsubsection*{Problem 4}
$a_n = 2\sqrt{2}\,a_{n-1}-4a_{n-2}$, $a_0=1$, $a_1=2$.

Characteristic equation: $r^2-2\sqrt{2}\,r+4=0$.
$r = \frac{2\sqrt{2}\pm\sqrt{8-16}}{2} = \frac{2\sqrt{2}\pm 2\sqrt{2}\,i}{2} = \sqrt{2}(1\pm i)$.

In polar form: $\rho=2$, $\theta=\pi/4$. So $a_n = 2^n(\alpha\cos(n\pi/4)+\beta\sin(n\pi/4))$.

$a_0 = \alpha = 1$. $a_1 = 2(\cos(\pi/4)+\beta\sin(\pi/4)) = 2\cdot\frac{\sqrt{2}}{2}(1+\beta) = \sqrt{2}(1+\beta)=2 \Rightarrow \beta = \sqrt{2}-1$.

$\boxed{a_n = 2^n\!\left(\cos\frac{n\pi}{4} + (\sqrt{2}-1)\sin\frac{n\pi}{4}\right)}$

\subsubsection*{Problem 5}
$a_n=2a_{n-1}+a_{n-2}$ (not standard Fibonacci), $a_0=0$, $a_1=1$.

Characteristic equation: $r^2-2r-1=0 \Rightarrow r=1\pm\sqrt{2}$.
General solution: $a_n=\alpha(1+\sqrt{2})^n+\beta(1-\sqrt{2})^n$.
$a_0=\alpha+\beta=0$, so $\beta=-\alpha$.
$a_1=\alpha(1+\sqrt{2})-\alpha(1-\sqrt{2})=2\alpha\sqrt{2}=1$, so $\alpha=\frac{1}{2\sqrt{2}}$.

$\boxed{a_n = \frac{(1+\sqrt{2})^n - (1-\sqrt{2})^n}{2\sqrt{2}}}$

\subsubsection*{Problem 6}
Binary strings of length 10 containing (at least) two consecutive zeros.

By complementary counting: let $b_n$ = strings of length $n$ with no ``00''.
$b_n = b_{n-1}+b_{n-2}$ (Fibonacci-like), $b_1=2$, $b_2=3$.
$b_{10}=144$. Total strings: $2^{10}=1024$.

$\boxed{1024-144=880}$

\subsubsection*{Problem 7}
Words of length $n$ from $\{a,b,c\}$ such that $a$ and $b$ are \emph{never} adjacent. (Interpreting ``$a$ and $b$ are adjacent'' as the constraint to count or avoid.)

\textit{Note:} The problem likely asks for words of a given length (say $n$) where $a$ is never adjacent to $b$. We set up a recurrence.

Let $f(n)$ = number of words of length $n$ where $a$ and $b$ are never adjacent. At each position, the letter is $a$, $b$, or $c$. If the current letter is $c$, the next can be anything valid; if $a$, the next cannot be $b$ and vice versa.

Let $A_n, B_n, C_n$ = words of length $n$ ending in $a, b, c$ respectively with no adjacent $a,b$.
$A_n = A_{n-1}+C_{n-1}$ (previous letter is $a$ or $c$).
$B_n = B_{n-1}+C_{n-1}$.
$C_n = A_{n-1}+B_{n-1}+C_{n-1}$.
$A_1=B_1=C_1=1$, so $f(1)=3$.

For general $n$, $f(n) = A_n+B_n+C_n$. The values grow; for $n=10$:

$n=1$: $A=1,B=1,C=1,f=3$.
$n=2$: $A=2,B=2,C=3,f=7$.
$n=3$: $A=5,B=5,C=7,f=17$.
$n=4$: $A=12,B=12,C=17,f=41$.
$n=5$: $A=29,B=29,C=41,f=99$.

The closed form satisfies $f(n)=2f(n-1)+f(n-2)$ with $f(1)=3, f(2)=7$, giving

\[
\boxed{f(n) = \frac{3+\sqrt{2}}{2\sqrt{2}}\,(1+\sqrt{2})^n + \frac{\sqrt{2}-3}{2\sqrt{2}}\,(1-\sqrt{2})^n}
\]

(If the problem intends a specific length, substitute accordingly; e.g.\ for length 10, $f(10)=8119$.)
